% SPDX-License-Identifier: Apache-2.0
% covers: Dm
\datasheet{Dm}{The RISC-V debug module on the bus: halt, resume and status, and the general registers, \code{dpc} and \code{dcsr} read and written by the abstract command.}

\dsfacts{\code{cpu/vreteno/src/debug.rs}}{\code{vreteno32::debug::Dm}}%
{\code{//docs:vreteno}}

\subsection*{Function}

The RISC-V debug specification puts a debugger's reach into a core
behind a debug module, a small address space of 32-bit registers a
transport carries from the cable. \code{Dm} is that module as far as
a debugger on the bus needs it, with the address space on the bus:
an AXI-Lite peripheral whose registers are the specification's, at
their numbers, each a word at four times its number. Through it the
JTAG host on the board's bus halts the core from the hardware manager,
reads its registers and resumes it (issue 154), and a transport of its
own reaches the same registers the same way when there is one.

The core's side is three lines and a word: a halt request and a
resume request in, debug mode out, and the register access, a number
in the specification's space, the word to write and whether to,
answered by the word read. The core lends the module its register
file's first read port and its write port while it is halted, and its
CSR read takes the module's number then.

\subsection*{Parameters}

None. One hart, so \code{hartsel} is not decoded; 32-bit registers,
so the command's \code{aarsize} must be 2.

\subsection*{Register map}

Offsets from the module's base, which on the board is
\code{0x1000\_0000}, a router port of its own; the part selects a
register by bits 8 to 2 of the address, the specification's number,
and looks at no other bits. \code{debug::at(number)} gives a
register's address on the board.

\begin{center}\footnotesize
\begin{tabular}{@{}lllp{0.45\textwidth}@{}}
\toprule
Offset & Number & Register & Access \\
\midrule
\code{0x10} & \code{0x04} & \code{data0} & Read and write, the
abstract command's one data word: what a read brought back, or what
a write sends. Written while a command is in flight it is the busy
error and unchanged. \\
\code{0x40} & \code{0x10} & \code{dmcontrol} & Read and write.
\code{haltreq} bit 31, a level; \code{resumereq} bit 30, an action,
which asks if the core is halted as the write lands and clears the
acknowledgement, and does nothing as a zero; \code{dmactive} bit 0,
which clear holds the module at nothing. \\
\code{0x44} & \code{0x11} & \code{dmstatus} & Read only.
\code{allhalted} and \code{anyhalted} bits 9 and 8, \code{allrunning}
and \code{anyrunning} bits 11 and 10, \code{allresumeack} and
\code{anyresumeack} bits 17 and 16, \code{authenticated} bit 7, and
the version, 2 for the 0.13 specification, in bits 3 to 0. \\
\code{0x48} & \code{0x12} & \code{hartinfo} & Reads as zero: no
scratch registers and no data window. \\
\code{0x58} & \code{0x16} & \code{abstractcs} & \code{busy} bit 12,
\code{cmderr} bits 10 to 8, cleared by writing ones, and
\code{datacount} 1 in bits 3 to 0. \\
\code{0x5c} & \code{0x17} & \code{command} & Write only, the
access-register command: \code{cmdtype} 0, \code{aarsize} 2,
\code{transfer} bit 17, \code{write} bit 16, \code{regno} bits 15 to
0, \code{0x1000} to \code{0x101f} a general register and below that
a CSR. \code{debug::access(regno, write)} builds one. \\
\code{0x100} & \code{0x40} & \code{haltsum0} & Read only, bit 0 the
one hart, halted. \\
\bottomrule
\end{tabular}
\end{center}

\dstables{Dm}

\subsection*{Behaviour}

\begin{itemize}
\item \code{haltreq} and \code{resumereq} go to the core as written
and held, and \code{halted} from the core is what \code{dmstatus}
reports. A resume request drops and is acknowledged when the core has
left debug mode; a halt request stays as written, and a debugger
clears it before it resumes.
\item A command is accepted when the module is active, nothing is in
flight and no error stands, and refused otherwise as the
specification says: \code{cmderr} 1 while one is in flight, 2 for a
command the module does not do, 4 while the core runs; under a
standing error a command is ignored. A command takes two cycles, one
for the number to reach the core and one for the word to come back
into \code{data0} or go in from it.
\item The module reads any general register and any CSR the core's
CSR read knows, and writes the general registers, \code{dpc} and
\code{dcsr}; a write of another CSR is \code{cmderr} 2.
\item The arms of its one \code{with!} are ordered for what may fire
together: the core's completion of a resume before the debugger's
write, so a request written in the cycle the core resumes clears the
acknowledgement rather than losing to it (issue 439).
\item A read is answered on \code{bus\_r} in the cycle its address beat is
taken, with \code{OKAY}; a write needs its address and data beats and
room on \code{bus\_b} in one cycle, and is answered \code{OKAY} there.
Inactive, the module holds nothing and answers zeros.
\end{itemize}

\subsection*{Verification}

\begin{itemize}
\item Thirteen tests in \code{cpu/vreteno/src/debug.rs}, in
\code{//cpu/vreteno:isa\_test}, play the core: they set halted and
the word read, drive the module over AXI-Lite, and read its lines
back. Among them:
\begin{itemize}
\item \code{a\_register\_is\_read\_and\_written\_by\_the\_abstract\_command};
\item \code{an\_unsupported\_command\_or\_a\_running\_core\_is}\-\code{\_an\_error\_that\_sticks};
\item \code{a\_resume\_request\_written\_as\_the\_core\_resumes}\-\code{\_clears\_the\_ack};
\item \code{data0\_written\_while\_a\_command\_is\_in\_flight}\-\code{\_is\_a\_busy\_error}.
\end{itemize}
The last two were constructed to fail against the arms as they were
before issue 439.
\item \code{the\_debug\_module\_halts\_reads\_writes\_and\_resumes\_the\_core},
in \code{//cpu/vreteno:board\_test}, puts a debugger on the board's
JTAG pins, a master model following a plan of writes, reads and
waits, against a program counting to three thousand: it halts the
core, reads the count and \code{dpc}, writes the count to its end,
resumes, and the program finishes at once.
\item \code{//cpu/vreteno:vreteno\_board\_dm\_probe} does the same from
Vivado's hardware manager through the JTAG host, against whatever
program runs; its transcript is owed to the document once the board
is reached again.
\end{itemize}

\subsection*{Used by}

\code{Dm} in the board's design, \code{Board}, as its field
\code{dmod} behind the bridge \code{pdm}, a \code{LiteBridge1} on the
router's seventh port at \code{0x1000\_0000}; its lines go to the
core's \code{haltreq}, \code{resumereq}, \code{dbg\_regno},
\code{dbg\_wdata} and \code{dbg\_we}, and take the core's \code{dbg}
and \code{dbg\_rdata}.

\subsection*{Limits}

One hart. The access-register command only: no program buffer, no
system bus access and no abstract access to memory, which the bus the
module sits on reaches directly. A CSR other than \code{dpc} and
\code{dcsr} is read but not written. No authentication, no reset of
the hart through \code{ndmreset}, no \code{havereset}. Not the transport:
the module is reached over the bus, and the debug transport module on
the cable is issue 154's third step. The module sits at
\code{0x1000\_0000} and not nearer the interrupt controller, whose
64~MiB window from \code{0x0c00\_0000} the router would match first.
